\documentclass{article}
\usepackage{geometry}
%\geometry{landscape}                   	
\geometry{width=520pt,columnsep=30pt}
\usepackage[utf8]{inputenc}
\usepackage{amssymb,amsmath}
\usepackage{hyperref}
\usepackage[all]{xypic}
\newtheorem{theorem}{Th}
\newtheorem{definition}{Def}

\def\AR#1{
\[
\begin{array}{llllllll}
#1
\end{array}
\]
}

\def\EQ#1{
\[
\begin{equation}
#1
\end{equation}
\]
}

\title{DN}
\author{theotimez}
\date{May 2020}


\begin{document}

\maketitle

\section{Behaviour of a clearing functor
under node fusion}

Consider a DN $(e,L)$, with $L_i$
a clearing vector.


Suppose we fuse nodes $a$, $b$ in $L$.
The fusion is described by
$f: I \to I\setminus\{a,b\}+\{[ab]\}:=I'$.


\subsection{fusion of solvables} 
Suppose both are solvable:
\AR{
e_a +
\sum_{j\neq a,b} \pi_{ja} L_j
+ \pi_{ba} L_b
&\geq&
\bar L_{a}
&
\hbox{because $a$ solvable}
\\
e_b +
\sum_{j\neq a,b} \pi_{jb} L_j
+ \pi_{ab} L_a
&\geq&
\bar L_{b}
&
\hbox{because $b$ solvable}
\\
e_a + e_b +
\sum_{j\neq a,b} (\pi_{ja}+\pi_{jb}) L_j
+ \pi_{ba} L_b
+ \pi_{ab} L_a
&\geq&
\bar L_{a}
+\bar L_{b}
\\
e_{[ab]} - \kappa(a,b) +
\sum_{j\neq a,b} (\pi_{ja}+\pi_{jb}) L_j
&\geq&
(\bar L_{a} - \pi_{ba} L_b)
+
(\bar L_{b} - \pi_{ab} L_a)
\\
&=&
\bar L_{a} - L_{ba}
+
\bar L_{b} - L_{ab}
\\
&=&
\bar L_{a,b}
}
donc les noeuds fusionnés restent L-solvables
dès que 
$\kappa(a,b)=0$.



Suppose $j\neq a, b$:
\AR{
e_j +
\sum_{k\neq a,b,j} 
\pi_{kj} L_j
+
\pi_{aj} \bar L_a
+
\pi_{bj} \bar L_b
&=&
e_j +
\sum_{k\neq a,b,j} 
\pi_{kj} L_j
+
L_{aj} +
L_{bj} 
\\
&=&
e_j +
\sum_{k\neq a,b,j} 
\pi_{kj} L_j
+
L_{[ab]j}
\\
&=&
e_j +
\sum_{k\neq a,b,j} 
\pi_{kj} L_j
+
\pi_{[ab]j} \bar L_{[ab]}
}


Donc si $L$ clear $(e,L)$,
alors $C(f)(L)$ clear 
$C(f)(e,L)$.

attention: on suppose les co\^uts de fusion nuls; a integrer

\subsection{fusion of unsolvables} 
Considérons le cas de noeuds insolvables.
On écrit le `revenu linéaire' $\phi_1(L)$ du noeud $j$ associé à un vecteur de paiement $L$ (rk: cette écriture convient aussi au cas $j=a$ ou $b$ car on a posé $\pi_{xx}=0$; rk2: puisqu'on suppose  $a$, $b$ insolvables, nécessairement $\bar L_a,\bar L_b>0$):
\AR{
e_j +
\sum_{k\neq a,b,j} 
\pi_{kj} L_k
+
\pi_{aj} L_a
+
\pi_{bj} L_b
&=&
e_j +
\sum_{k\neq a,b,j} 
\pi_{kj} L_k
+
\frac{L_{aj}}{
(\sum_{k\neq b} L_{ak})
+ L_{ab}
} L_a
+
\frac{L_{bj}}{
(\sum_{k\neq a} L_{bk})
+ L_{ba}
} 
L_b
\\
e_a +
\sum_{k\neq a,b} 
\pi_{ka} L_k
+
\pi_{ba} L_b
&=&
e_a +
\sum_{k\neq a,b} 
\pi_{ka} L_k
+
\frac{L_{ba}}{
(\sum_{k\neq a} L_{bk})
+ L_{ba}
} 
L_b
}
On écrit maintenant le revenu de $j$
sous le paiement $C(f)(L)$ après fusion pour $j\neq a,b$:
\AR{
e_j +
\sum_{k\neq a,b,j} 
\pi_{kj} L_k
+
\pi_{[ab]j} L_{[ab]}
}

On a:
\AR{
\pi_{[ab]j} &=&
\frac
{L_{[ab]j}}
{\sum_k L_{[ab]k}}
&=&
\frac
{L_{aj}+L_{bj}}
{
\sum_k L_{ak} + L_{bk} -
(L_{ab}+L_{ba})
}
}
donc le revenu de $j$ -
$\phi_1(C(f)(L))$:
\AR{
e_j +
\sum_{k\neq a,b,j} 
\pi_{kj} L_k
+
\frac
{L_{aj}+L_{bj}}
{
\sum_k L_{ak} + L_{bk} -
L_{ab}- L_{ba}
}
(L_a+L_b)
}

le revenu de $[ab]$ après fusion (no fusion cost):
\AR{
e_a + e_b + 
\sum_{k\neq a,b} 
(\pi_{ka}+\pi_{kb}) L_k
}

RK: $C(f)$ envoie l'hyperrectangle $[0,\bar L]$
dans un hyper-rectangle qui 
contient strictement, si
$L_{ab}+L_{ba}>0$, celui du système fusionné.

Comparons les revenus linéaires de $j\neq a,b$ dans les deux scenarios:

\AR{
\frac{L_{aj}}{
(\sum_{k\neq b} L_{ak})
+ L_{ab}
} L_a
+
\frac{L_{bj}}{
(\sum_{k\neq a} L_{bk})
+ L_{ba}
} 
L_b
&\leq\geq?&
\frac
{L_{aj}+L_{bj}}
{
\sum_k L_{ak} + L_{bk} -
L_{ab}- L_{ba}
}
(L_a+L_b)
\\
%e_a + e_b + 
%\sum_{k\neq a,b} 
%(\pi_{ka}+\pi_{kb}) L_k
%+
%\frac{L_{ba}}{
%(\sum_{k\neq a} L_{bk})
%+ L_{ba}}
%+
%\frac{L_{ab}}{
%(\sum_{k\neq b} L_{ak})
%+ L_{ab}}
\pi_{ba} L_b +
\pi_{ab} L_a
&\leq\geq?&
0
%e_a + e_b + 
%\sum_{k\neq a,b} 
%(\pi_{ka}+\pi_{kb}) L_k
}

\AR{
\xymatrix@C=50pt{
1
\ar[r]^{\bar L}
\ar[d]^{=}
&
\mathbb R_+^I 
\ar[r]^{\phi}
\ar[d]_{C(f)}
&
\mathbb R_+^I
\ar[r]^{max(\_,\bar  L)}
\ar[d]^{C(f)}
&
\mathbb R_+^I
\ar[d]^{C(f)}
\\
1
\ar[r]^{C(f)(\bar  L)}
&
\mathbb R_+^{I'} 
\ar[r]^{\phi'_1}
&
\mathbb R_+^{I'}
\ar[r]^{max(\_,C(f)(\bar  L))}
&
\mathbb R_+^{I'}
}
}

idealement $\phi'\circ f \geq f\circ\phi$, ce qui voudrait dire que
les orbites basses sont Pareto-improving. Mais ça ne peut pas etre vrai si $a$, $b$ solvables; donc on va essayer de démontrer cette inégalité quand au moins l'un des deux est insolvable.

Reprenons l'approche par plus petit point fixe car les $\phi^n(0)$ 
représentent des vrais paiements et essayons de raisonner directement
sur la valeur $v(x) = \sum_n \phi^{n+1}(0)_x-\phi^n(0)_x$?


\subsection{Questions}

Question liée à la fragilité d'un réseau

Étant donné un DN solvable, et un noeud $i$,
quels sont les e extremaux pour $i$?
$i$ est solvable dans $(L,e)$ mais pas
pour $(L,e'<e)$.

quelle regle pour organiser la fusion?

enchere sur les banques sauveteurs?

est-ce qu'une fusion d'insolvables 
peut devenir solvable??

\section{next}

lire \S4 de Rogers.\\

Soit $N = \lbrace 1;...; n\rbrace$ un ensemble de banque et $(L; e; \alpha; \beta)$ un système financier correspondant.\\
Supposons les banques $1$ et $n$ $0$-insolvables, c'est à dire :
\begin{center}
$\bar{L}_{1} > e_{1} + \sum_{j=1}^{n} \bar{L}_{j} \pi_{j1} = e_{1} + \sum_{j=1}^{n} L_{j1}$\\
$\bar{L}_{n} > e_{n} + \sum_{j=1}^{n} \bar{L}_{j} \pi_{jn} = e_{n} + \sum_{j=1}^{n} L_{jn}$
\end{center}
Donc :
\begin{center}
 $\bar{L}_{1} + \bar{L}_{n} > e_{1} + e_{n} + \sum_{j=1}^{n} L_{j1} + \sum_{j=1}^{n} L_{jn}$\\
$\bar{L}_{1} + \bar{L}_{n} - L_{n1} - L_{1n} > e_{1} + e_{n} + \sum_{j=1}^{n-1} L_{j1} + \sum_{j=2}^{n} L_{jn}$\\
$\bar{L}_{\left[ 1;n\right]} > e_{\left[ 1;n\right]} + \sum_{j=2}^{n-1} L_{j1} + \sum_{j=2}^{n-1} L_{jn}$ \\
$\bar{L}_{\left[ 1;n\right]} > e_{\left[ 1;n\right]} + \sum_{j=2}^{n-1} L_{j\left[ 1;n\right]}$
\end{center}
Par conséquent la fusion de deux nœuds $0$-insolvables est aussi insolvable.

\subsection{Questions}

comparer les deux manques à payer?

Q2 - suppose $L_{1n}=L_{n1}\to\infty$;
le clearing va etre $\bar L$ pour les $i\neq 1,n$; et tout la dette  interne va etre $1\to n$ et $n\to 1$ (on est des puits);
$\pi_{1n}=1=\pi_{n1}$;
si tous les autres noeuds ont de quoi payer sans la contribution de $1$ et $n$,
alors on a un point fixe en deux coups
et c'est un puits.

scenario a 3 -
Supposons 3 noeuds: 
1 tres solvable, 
2 insolvables qui forment un ``puits''
en se pretant beaucoup.

scenario a 2 - Pire??: est-ce qu'un noeud insolvable peut aller contracter un accord secret avec un noeud solvable pour former 
une grosse dette fictive et obtenir un effet ``puits'' la somme (officieuse) des deux noeuds dérive plus de thune qu'avant. Remarque - la collusion peut backfire sur la banque solvable si son ``pompage'' conduit à la banqueroute de certains de ses créanciers.


Plus generalement, quel est la sensibilité du vecteur de $L^\star$ aux variations conservatives de la matrice $L_{ij}$, $L \to L + S$ ou $S$
symetrique.




\subsection{4 possibilities}
Soit $L$ un système financier à $n$ banques toutes solvables.

Alors soit i une banque on a 4 possibilités :

(1) $e_{i} \geq \bar{L}_{i} +\alpha$ et $\sum L_{ji} + \alpha \geq \bar{L_{i}} + \alpha$

(2)  $e_{i} \geq \bar{L}_{i}$ et $\sum L_{ji} < \bar{L_{i}}$

(3) $e_{i} < \bar{L}_{i}$ et $\sum L_{ji} \geq \bar{L_{i}}$

(4)  $e_{i} < \bar{L}_{i}$ et $\sum L_{ji} < \bar{L_{i}}$

Dans le cas (1), i ne peut pas devenir insolvable par seule baisse d'assets donc i ne peut être 0-insolvable.

 Si la baisse est suffisamment forte alors i se retrouve dans le cas (3).
 
 Sinon i reste dans le cas (1) et est donc solvable. On dira que i est 0-solvable (on sait dès la baisse d'assets que i est et restera solvable).
 
Dans le cas (2), i peut devenir insolvable par une baisse d'assets et est alors 0-insolvable.

Si i reste solvable alors :\\
Si la baisse est suffisamment forte, i se retrouve dans le cas (4).\\
Sinon i reste dans le cas (2) et est alors 0-solvable.

Dans le cas (3), i ne peut pas devenir insolvable par seule baisse d'assets donc n'est jamais 0-insolvable mais elle peut se retrouver 1-insolvable par contagion.

Dans le cas (4), i peut devenir insolvable par baisse d'asset ou par contagion.


\subsection{exemple ou les valeurs dépendent du point fixe choisi.}

A priori quand $$\alpha=\beta=1$$ ce n'est pas possible.

$(x\; x)$ est un clearing vector:
\AR{
x + e < \gamma
&\Rightarrow& 
\phi(x) = \beta x + \alpha e
\\
x  &=& 
\alpha e/(1 - \beta)
& \text{supposing } \beta<1
\\
x + e \geq \gamma 
&\Rightarrow& 
\phi(x) = \gamma
\\
x = \gamma
}
with eg $\gamma=2.2$.  

\begin{equation*}
\textbf{L} = 
\begin{pmatrix}
0 & 2.2 \\
2.2 & 0 \\
\end{pmatrix}
\end{equation*}

\begin{equation*}
\Pi = 
\begin{pmatrix}
0 & 1 \\
1 & 0 \\
\end{pmatrix}
\end{equation*}

\begin{equation*}
\textbf{e} = 
\begin{pmatrix}
1 \\
1 \\
\end{pmatrix}
\end{equation*}

\begin{equation*}
\alpha = 0.5   \
\
\beta = 0.5
\end{equation*}

D'après l'article 
$\begin{pmatrix}
2.2 \\
2.2 \\
\end{pmatrix}$ 
et $\begin{pmatrix}
1 \\
1 \\
\end{pmatrix}$
sont deux clearing vecteurs.

Or $\begin{pmatrix}
0 & 1 \\
1 & 0 \\
\end{pmatrix}$
 $\begin{pmatrix}
1 \\
1 \\
\end{pmatrix}$ +
$\begin{pmatrix}
1 \\
1 \\
\end{pmatrix}$ -
$\begin{pmatrix}
2.2 \\
2.2 \\
\end{pmatrix}$ = 
-$\begin{pmatrix}
0.2 \\
0.2 \\
\end{pmatrix}$

et 
\
$\begin{pmatrix}
0 & 1 \\
1 & 0 \\
\end{pmatrix}$
 $\begin{pmatrix}
2.2 \\
2.2 \\
\end{pmatrix}$ +
$\begin{pmatrix}
1 \\
1 \\
\end{pmatrix}$ -
$\begin{pmatrix}
2.2 \\
2.2 \\
\end{pmatrix}$ = 
$\begin{pmatrix}
1 \\
1 \\
\end{pmatrix}$

Par conséquent le théorème d'égalité n'est pas vrai pour $\alpha, \beta \neq 1$

\textbf{RK}: est ce que le plus grand point fixe conduit toujours à un meilleur payoff pour tous les joueurs (Pareto improving)? Oui, car
le revenu final d'une banque est une fonction croissante du vecteur de clearing.



\subsection{La fusion de deux $\mu$ -insolvables est $\mu$ -insolvable}

Si 1 et n sont deux banques $\mu$-insolvables alors:

$\sum_{j=1}^{n} \Lambda^{(\mu)}_{j} \pi_{j1} + e_{1} < \bar L_{1}$

$\sum_{j=1}^{n} \Lambda^{(\mu)}_{j} \pi_{jn} + e_{n} < \bar L_{n}$

Donc: 
$\sum_{j=1}^{n} (\Lambda^{(\mu)}_{j} \pi_{j1} + \Lambda^{(\mu)}_{j} \pi_{jn}) + e_{1}  + e_{n} < \bar L_{1} + \bar L_{n}$

On suppose $1$ et $n$, $(\mu-1)$-solvables par conséquent:
$\Lambda^{(\mu)}_{1}=\bar  L_1$ et
$\Lambda^{(\mu)}_{n}=\bar  L_n$; 
donc
$\Lambda^{(\mu)}_{n} \pi_{n1} = L_{n1}$ et 
$\Lambda^{(\mu)}_{1} \pi_{1n} = L_{1n}$.

Et donc le noeud fusionné doit la somme suivante au reste du monde (sous le vecteur $\Lambda^{(\mu)})$: 
$\sum_{j=2}^{n-1} (\Lambda^{(\mu)}_{j} \pi_{j1} + \Lambda^{(\mu)}_{j} \pi_{jn}) + e_{1}  + e_{n} < \bar L_{1} + \bar L_{n} - L_{n1} - L_{1n} = \bar L_{\left[1n \right]}$

($\Lambda^{(\mu)}_{j}$ décroît jusqu'au plus grand vecteur de clearing donc le membre de gauche est supérieur ou égal à $\Phi_{1}(L^{*})$).

Ce résultat indique que deux banque $\mu$-insolvables qui risque d'être contagionnées par l'insolvabilité des banques $<\mu$-insolvables ne peuvent stopper la propagation en fusionnant ensemble (ce qui leurs aurait permi de laisser les banques les plus insolvables sombrer).

\textbf{Question} : la fusion d'une banque $\mu$-insolvable et d'une $>\mu+1$-insolvable peut-elle être solvable? Si $\mu > 0$ cela revient à laisser les banques 0-insolvables mourir, si $\mu$=0 cela reviendrait à peut-être sauver aussi les banques $<\mu+1$-insolvables sans que celles ci fusionnent, ou alors à sauver des banques sélectivement suivant leur niveau d'insolvabilité..


paiements paradoxaux. Au début T, V ont 1 et la dette mutuelle est 2:
\AR{
%T\to V: VT2?
%\\
%V\to T: TV1?
%\\
T\to V: TV1!
\\
V\to T: VT2!
\\
T\to V: TV2!
\\
V\to T: VT1!
}


\subsection{Fusion d'une banque 0-insolvable avec une 1-insolvable}

Soit 1 la banque 0-insolvable donc 
$\sum_{j=1}^{n} L_{j1} + e_{1} < \bar L_{1}$.

Soit n la banque 1-insolvable donc 

$\sum_{j=1}^{n} \Lambda^{(\mu)}_{j} \pi_{jn} + e_{n} = \sum_{j \in S_{0}} L_{jn} + \sum_{j \in I_{0}} x_{j} \pi_{jn} + e_{n} < \bar L_{n}$

On a donc :

$e_{1} + e_{n} + \sum_{j \in S_{0}} (L_{j1} + L_{jn}) + \sum_{j \in I_{0}} (L_{j1} + x_{j} \pi_{jn})  < \bar L_{1} +\bar L_{n}$ 

$e_{1} + e_{n} + \sum_{j \in S_{0}\setminus \lbrace n \rbrace} (L_{j1} + L_{jn}) + \sum_{j \in I_{0}} (L_{j1} + x_{j} \pi_{jn}) - L_{1n} < \bar L_{\left[1n \right]} = \bar L_{1} +\bar L_{n} - L_{1n} - L_{n1}$

Or la richesse de la fusion de 1 et n est : 

$e_{1} - \kappa_{1} + e_{n} - \kappa_{n} + \sum_{j \in S_{0}\setminus \lbrace n \rbrace} (L_{j1} + L_{jn}) + \sum_{j \in I_{0}\setminus \lbrace 1 \rbrace} x_{j}(\pi_{j1} + \pi_{jn})$

On souhaite savoir sous quelles conditions on a :

$e_{1} - \kappa_{1} + e_{n} - \kappa_{n} + \sum_{j \in S_{0}\setminus \lbrace n \rbrace} (L_{j1} + L_{jn}) + \sum_{j \in I_{0}\setminus \lbrace 1 \rbrace} x_{j}(\pi_{j1} + \pi_{jn}) \geq \bar L_{\left[1n \right]}$
 
Il faut donc : 

$- \kappa_{1} - \kappa_{n} + \sum_{j \in I_{0}\setminus \lbrace 1 \rbrace} x_{j}(\pi_{j1} + \pi_{jn}) > \sum_{j \in I_{0}} (L_{j1} + x_{j} \pi_{jn}) - L_{1n}$

$- \kappa_{1} - \kappa_{n} + \sum_{j \in I_{0}\setminus \lbrace 1 \rbrace} x_{j} \pi_{j1} > \sum_{j \in I_{0}} L_{j1} + x_{1} \pi_{1n} - L_{1n}$

$L_{1n} - x_{1} \pi_{1n}> \kappa_{1} + \kappa_{n} + \sum_{j \in I_{0}\setminus \lbrace 1 \rbrace} (L_{j1} - x_{j} \pi_{j1})  $ \textit{(condition nécessaire)}

(à gauche la différence entre ce que doit 1 à n et ce que paye 1 à n. A droite la différence entre ce que 1 devait recevoir et ce que 1 reçoit vraiment, plus les coûts de fusion).


La fusion est solvable si et seulement si :

$e_{1} - \kappa_{1} + e_{n} - \kappa_{n} + \sum_{j \in S_{0}\setminus \lbrace n \rbrace} (L_{j1} + L_{jn}) + \sum_{j \in I_{0}\setminus \lbrace 1 \rbrace} x_{j}(\pi_{j1} + \pi_{jn}) \geq \bar L_{\left[1n \right]}$

$e_{1} - \kappa_{1} + e_{n} - \kappa_{n} + \sum_{j \in S_{0}} (L_{j1} + L_{jn}) + \sum_{j \in I_{0}\setminus \lbrace 1 \rbrace} x_{j}(\pi_{j1} + \pi_{jn}) \geq \bar L_{1} + \bar L_{n}$

$ e_{n} - \kappa_{n} + \sum_{j \in S_{0}} L_{jn} + \sum_{j \in I_{0}\setminus \lbrace 1 \rbrace} x_{j}\pi_{jn} - \bar L_{n} \geq \bar L_{1}  - (\sum_{j \in S_{0}} L_{j1} + \sum_{j \in I_{0}\setminus \lbrace 1 \rbrace} x_{j}\pi_{j1} + e_{1} - \kappa_{1})$

$ e_{n} - \kappa_{n} + \sum_{j=1}^{n} L_{jn} - (\sum_{j \in I_{0}\setminus \lbrace 1 \rbrace} L_{jn} - x_{j}\pi_{jn}) - \bar L_{n} \geq \bar L_{1}  - (\sum_{j=1}^{n} L_{j1} - \sum_{j \in I_{0}\setminus \lbrace 1 \rbrace} (L_{j1} - x_{j}\pi_{j1})  + e_{1} - \kappa_{1})$

$ e_{n}  + \sum_{j=1}^{n} L_{jn} - \bar L_{n}- \kappa_{n} - (\sum_{j \in I_{0}\setminus \lbrace 1 \rbrace} L_{jn} - x_{j}\pi_{jn})  \geq \bar L_{1}  - (\sum_{j=1}^{n} L_{j1} + e_{1} - \sum_{j \in I_{0}\setminus \lbrace 1 \rbrace} (L_{j1} - x_{j}\pi_{j1})   - \kappa_{1})$

Question? Une fusion partielle de nœuds insolvables (à des niveaux différents) peut-elle rendre ces banques définitivement solvable? Sous quelles conditions? Comment se comporte le niveau d'insolvabilité? Comment la fusion affecte le niveau d'insolvabilité : la fusion devient solvable au niveau 0 mais si une nouvelle baisse se produit la banque fusionnée serait-elle au mieux 1-insolvable? La fusion peut-elle changer le type de banque (parmi les 4) et donc modifier la sensibilité aux risques?

\subsection{Cas général}
Soient $(L, e, \alpha, \beta)$ et $(L', e', \alpha, \beta)$ deux systèmes financiers à $n$ banques et $n-1$ banques respectivement. $(L', e', \alpha, \beta)$ est le système $(L, e, \alpha, \beta)$ après fusion des banques $1$ et $n$, notée $\left[1n \right]$ et correspondant à la banque $1$ dans le système $L'$.

Supposons $1$ et $n$ de $(L, e, \alpha, \beta)$, 0-insolvables. Alors : 

\begin{center}
$\bar{L}_{1} > e_{1} + \sum_{j=1}^{n} \bar{L}_{j} \pi_{j1} = e_{1} + \sum_{j=1}^{n} L_{j1}$

$\bar{L}_{n} > e_{n} + \sum_{j=1}^{n} \bar{L}_{j} \pi_{jn} = e_{n} + \sum_{j=1}^{n} L_{jn}$
\end{center}


Donc la richesse de la banque $\left[1n \right]$ dans le système $(L', e', \alpha, \beta)$ est :
\begin{center}
 $e'_{1} + \sum_{j=1}^{n-1} L'_{j1} = e_{1} - \kappa_{1} + e_{n} - \kappa_{n} + \sum_{j=2}^{n-1} (L_{j1} + L_{jn}) <  \bar{L}_{1} + \bar{L}_{n} - L_{n1} - L_{1n} = \bar{L}'_{1}$
 \end{center}
Donc $e'_{1} + \sum_{j=1}^{n-1} L'_{j1} < \bar{L}'_{1}$

Par conséquent 1 est 0-insolvable dans le système $(L', e', \alpha, \beta)$.

Supposons $1$ et $n$ de $(L, e, \alpha, \beta)$, $(\mu -1)$-solvables et $\mu$-insolvables (pour $\mu > 0$). Alors : 
\begin{center}
$\sum_{j=1}^{n} \Lambda^{(\mu)}_{j} \pi_{j1} + e_{1} < \bar L_{1}$

$\sum_{j=1}^{n} \Lambda^{(\mu)}_{j} \pi_{jn} + e_{n} < \bar L_{n}$
\end{center}
et 
\begin{center}
$\sum_{j=1}^{n} \Lambda^{(\mu-1)}_{j} \pi_{j1} + e_{1} \geq \bar L_{1}$

$\sum_{j=1}^{n} \Lambda^{(\mu-1)}_{j} \pi_{jn} + e_{n} \geq \bar L_{n}$
\end{center}

Donc la richesse de la banque $\left[1n \right]$ dans le système $(L', e', \alpha, \beta)$ à l'étape $\theta$ est :
\begin{center}
 $e'_{1} + \sum_{j=1}^{n-1} \Lambda'^{(\theta)}_{j} \pi'_{j1} = e_{1} - \kappa_{1} + e_{n} - \kappa_{n} + \sum_{j=2}^{n-1} \Lambda^{(\theta)}_{j} (\pi_{j1} + \pi_{jn})$
 \end{center}

On remarque pour $\theta = \mu - 1$
\begin{center}
$e'_{1} + \sum_{j=1}^{n-1} \Lambda'^{(\mu-1)}_{j} \pi'_{j1} = e_{1} - \kappa_{1} + e_{n} - \kappa_{n} + \sum_{j=2}^{n-1} \Lambda^{(\mu-1)}_{j} (\pi_{j1} + \pi_{jn})$

et

$e_{1} - \kappa_{1} + e_{n} - \kappa_{n} + \sum_{j=2}^{n-1} \Lambda^{(\mu-1)}_{j} (\pi_{j1} + \pi_{jn}) \geq \bar L_{1} + \bar L_{n} - L_{1n} - L_{n1} - (\kappa_{1} + \kappa_{n}) = \bar L'_{1} - (\kappa_{1} + \kappa_{n})$

donc

$e'_{1} + \sum_{j=1}^{n-1} \Lambda'^{(\mu-1)}_{j} \pi'_{j1} \geq \bar L'_{1} - (\kappa_{1} + \kappa_{n})$
\end{center}

Il n'est donc pas garanti que la fusion de deux $\mu-1$-solvables le soit.

Supposons que ce soit le cas, c'est à dire :

$e'_{1} + \sum_{j=1}^{n-1} \Lambda'^{(\mu-1)}_{j} \pi'_{j1} \geq \bar L'_{1}$

Alors on a : 
\begin{center}
$e'_{1} + \sum_{j=1}^{n-1} \Lambda'^{(\mu)}_{j} \pi'_{j1} = e_{1} - \kappa_{1} + e_{n} - \kappa_{n} + \sum_{j=2}^{n-1} \Lambda^{(\mu)}_{j} (\pi_{j1} + \pi_{jn}) \leq e_{1} + e_{n} + \sum_{j=2}^{n-1} \Lambda^{(\mu)}_{j} (\pi_{j1} + \pi_{jn}) < \bar{L}_{1} + \bar{L}_{n} - L_{1n} - L_{n1} = \bar L'_{1} $ 
\end{center}
Donc $e'_{1} + \sum_{j=1}^{n-1} \Lambda'^{(\mu)}_{j} \pi'_{j1} < \bar L'_{1}$, et la banque fusionnée est donc $\mu$-insolvable dans le système $(L', e', \alpha, \beta)$.

La condition que la fusion soit aussi solvable que le sont les banques signifie que la somme des richesses des banques 1 et n à l'étape $\mu$ est supérieure à la somme de leurs dettes et des coûts de fusion :
\begin{center}
La fusion de 1 et n $\mu$-solvables est $\mu$-solvable ssi 

$(\sum_{j=1}^{n} \Lambda^{(\mu)}_{j} \pi_{j1} + e_{1}) + (\sum_{j=1}^{n} \Lambda^{(\mu)}_{j} \pi_{jn} + e_{n}) \geq \bar L_{1} + \bar L_{n} + (\kappa_{1} + \kappa_{n})$
\end{center}

Or $\Lambda^{(\mu)}$ décroît en fonction de $\mu$ par conséquent plus les banques deviennent insolvables tard et plus les coûts de fusion doivent être faibles.

\subsection{système à deux banques}

Soient
\begin{equation*}
\textbf{L} = 
\begin{pmatrix}
0 & L_{12} \\
L_{21} & 0 \\
\end{pmatrix}
\end{equation*}

\begin{equation*}
\Pi = 
\begin{pmatrix}
0 & 1 \\
1 & 0 \\
\end{pmatrix}
\end{equation*}

\begin{equation*}
\textbf{e} = 
\begin{pmatrix}
e_{1} \\
e_{n} \\
\end{pmatrix}
\end{equation*}

\begin{equation*}
\alpha    \
\
\beta
\end{equation*}

On a alors 4 possibilités :

Le système est solvable et donc $L' = (\bar{L}_{1} ; \bar{L}_{2})$

La banque 1 est solvable et la banque 2 est insolvable donc 
\begin{center}
$e_{1} + L'_{2} \geq \bar{L}_{1} = L_{21}$ donc $L'_{1} = L_{21}$

$e_{2} + L_{21} < \bar{L}_{2}$ donc $L'_{2} = \alpha e_{2} + \beta L_{21}$
\end{center}

La banque 2 est solvable et la banque 1 est insolvable donc 
\begin{center}
$e_{2} + L'_{1} \geq \bar{L}_{2} = L_{12}$ donc $L'_{2} = L_{12}$

$e_{1} + L_{12} < \bar{L}_{1}$ donc $L'_{1} = \alpha e_{1} + \beta L_{12}$
\end{center}

Aucune des banques n'est solvable donc
\AR{
L'_{1} = \alpha e_{1} + \beta L'_{2} 
\\
L'_{2} = \alpha e_{2} + \beta L'_{1}
}

Si les deux sont insolvables:
\AR{
e_{1} + L_{21} < L_{12}
\\
e_{2} + L_{12} < L_{21}
}
c'est impossible!

$\Longrightarrow$ Or $L'_{1} = \alpha \frac{\beta e_{2} + e_{1}}{1- \beta^{2}} $ et $L'_{2} = \alpha \frac{\beta e_{1} + e_{2}}{1 - \beta^{2}}$
est un point fixe de $\phi_{\alpha,\beta}$ à condition que:
\AR{
 e_1 +
 \alpha \frac{\beta e_{1} + e_{2}}{1 - \beta^{2}}
 < L_{12}
\\
 e_2 +
 \alpha \frac{\beta e_{2} + e_{1}}{1 - \beta^{2}}
 < L_{21}
 }

On a un exemple de vecteur de clearing fallacieux: au sens où il prédit qu'un noeud est insolvable alors que ce n'est pas forcément le cas (en choisissant un autre vecteur au-dessus dans l'ordre produit)

Q: Où va la richesse évaporée 
$L'_1+L'_2= \dfrac{\alpha(e_1+e_2)}{1-\beta}$ (diverge qd $\beta\to 1^-$!!)
si tout le monde a disparu?

si $\beta=1$
\AR{
x_1 = e_1 + L_{21} < L_{12} < L_{21} - e_2\\
x_2 = e_2 + L_{12} < L_{21}
}
conclusion il ne peut pas y avoir de double insolvabilité dans le cas $\beta=1$, $\alpha >0$.

\subsection{magouille prêt symètrique}
On considère un système $(L; e ; \alpha ; \beta)$ à $n$ banques avec une seule banque 0-insolvable : la banque 1, donc :

\AR{
\sum_{j=1}^{n} \Lambda^{(0)}_{j} \pi_{j1} + e_{1} <  \bar{L}_{1}
& \text{1 est 0-insolvable dans $L$ et donc dans $L'$}
\\
\Lambda^{(1)}_{1} = x_{1} = \alpha e_{1} + \beta \sum_{j=1}^{n} L_{j1}
}

On rajoute à L, une matrice de taille $n\times n$  dont tous les coefficients sont nuls exceptés $a_{1n}=S$ et $a_{n1}=S$, on note $L'$ le résultat (dette fictive $S$). ca ne peut pas rendre 1 solvable (voir ci-dessous). Par ailleurs 
$\pi_{1n}(S)\to_{S\to\infty} 1$,
donc la banque $n$ (si elle reste solvable)
peut .

$(L'; e ; \alpha ; \beta)$ est un système financier représentant où les banques 1 et n ont contracté une dette mutuelle.


Si $i \notin \left\lbrace 1; n \right\rbrace$ alors $\bar{L'}_{i} = \bar{L}_{i}$, sinon $\bar{L'}_{i} = \bar {L}_{i} + S$

Si $i;j \notin \left\lbrace 1; n \right\rbrace$ alors $L'_{ij} = L_{ij}$.


Une banque est solvable à l'étape 0 dans le système L, si et seulement si, elle l'est dans le système L'.


$\sum_{j=1}^{n} L_{ji} + e_{i} <  \bar{L}_{i}$ donc $i \notin \left\lbrace 1; n \right\rbrace$ c'est bon

sinon $\sum_{j=1}^{n} L_{ji} + S + e_{i} <  \bar{L}_{i} + S = \bar{L'}_{i}$
 
On fait pareil pour le cas supérieur ou égal.

\begin{center}
Si $j \notin \left\lbrace 1; n \right\rbrace$ : $\Lambda^{(0)}_{j} = \bar{L}_{j} = \bar{L'}_{j} = \Lambda'^{(0)}_{j} $

Sinon $\Lambda^{(0)}_{j} + S = \bar{L}_{j} + S = \bar{L'}_{j} = \Lambda'^{(0)}_{j} $
\end{center}

\begin{center}
Si $j \notin \left\lbrace 1; n \right\rbrace$ : $\Lambda^{(1)}_{j} = \bar{L}_{j} = \bar{L'}_{j} = \Lambda'^{(1)}_{j} $

Si $j=n$ $\Lambda^{(1)}_{n} + S = \bar{L}_{j} + S = \bar{L'}_{j} = \Lambda'^{(1)}_{j} $

Si $j=1$ $ \Lambda^{(1)}_{1} = x_{1} = \alpha e_{1} + \beta \sum_{j=1}^{n} L_{j1}$ 

$ \Lambda'^{(1)}_{1} = \alpha e_{1} + \beta ( \sum_{j=1}^{n-1} L_{j1} + \bar{L'}_{n} \pi'_{n1}) = \Lambda^{(1)}_{1} + \beta S $
\end{center}


Supposons que le système L soit solvable à l'étape 1 donc $\Lambda^{(1)}$ est le plus grand vecteur de clearing. 

Cela signifie donc qu'à l'étape 1, la seule banque insolvable pour $\Lambda^{(1)}$ est la 1.

Vérifions si dans le sytème L' la solvabilité des autres banques se maintient.
\begin{center}
Pour $i \neq n$ on a $ \sum_{j=1}^{n} \Lambda^{(1)}_{j} \pi_{ji} + e_{i} \geq \bar{L}_{i}$

donc $ \sum_{j=2}^{n-1} \Lambda^{(1)}_{j} \pi_{ji} + \Lambda^{(1)}_{1} \pi_{1i} + \Lambda^{(1)}_{n} \pi_{ni} \geq \bar{L}_{i}$ 
\end{center}
or $ \sum_{j=2}^{n-1} \Lambda^{(1)}_{j} \pi_{ji} =  \sum_{j=2}^{n-1} \Lambda'^{(1)}_{j} \pi'_{ji}$

$\Lambda^{(1)}_{n} \pi_{ni} = L_{ni} = \Lambda'^{(1)}_{n} \pi'_{ni}$

$\Lambda'^{(1)}_{1} \pi'_{1i} = (\Lambda^{(1)}_{1} + \beta S) \frac{L_{1j}}{\bar{L}_{1}+S}$

\begin{center}
Si $\beta \geq \frac{\Lambda^{(1)}_{1}}{\bar{L}_{1}}$

$\beta S \geq \frac{\Lambda^{(1)}_{1}}{\bar{L}_{1}} S$

$\Lambda'^{(1)}_{1} - \Lambda^{(1)}_{1} \geq \frac{\Lambda^{(1)}_{1}}{\bar{L}_{1}}S$

$\Lambda'^{(1)}_{1} \geq \Lambda^{(1)}_{1} (1 + \frac{S}{\bar{L}_{1}})$

$\Lambda'^{(1)}_{1} \geq \Lambda^{(1)}_{1} \frac{\bar{L}_{1} + S}{\bar{L}_{1}}$

$\Lambda'^{(1)}_{1} \frac{L_{1i}}{\bar{L}_{1}+S} \geq \Lambda^{(1)}_{1} \frac{L_{1i}}{\bar{L}_{1}}$

$\Lambda'^{(1)}_{1} \pi'_{1i} \geq \Lambda^{(1)}_{1} \pi_{1i}$
\end{center}

Donc $\beta \geq \frac{\Lambda^{(1)}_{1}}{\bar{L}_{1}} \Rightarrow$ i solvable dans L'.

Question : On a une condition suffisante mais reste à trouver une nécessaire.

\begin{center}
Pour $i = n$ on a $ \sum_{j=1}^{n} \Lambda^{(1)}_{j} \pi_{jn} + e_{n} \geq \bar{L}_{n}$

donc $ \sum_{j=2}^{n-1} \Lambda'^{(1)}_{j} \pi'_{ji} + \Lambda^{(1)}_{1} \pi_{1n} + S \geq \bar{L'}_{n}$ 
\end{center}

Or $1 \geq \beta \pi'_{1n}$

$\frac{S}{\pi'_{1n}} \geq \beta S$

$\frac{S}{\pi'_{1n}} \geq \Lambda'^{(1)}_{1} - \Lambda^{(1)}_{1}$

$S \geq \Lambda'^{(1)}_{1} \pi'_{1n} - \Lambda^{(1)}_{1} \pi'_{1n} \geq \Lambda'^{(1)}_{1} \pi'_{1n} - \Lambda^{(1)}_{1} \pi_{1n}$

Donc pas de garantie que n reste solvable (peut-être même impossible?)

\subsection{magouille à 3}
Consider a system with 3 nodes $1$, $2$ and $n$ with only non zero coefficients $L_{12}$, $L_{21}$, $L_{1n}$; and apply the magouille transform; then profitability for $n$ reads (assuming $n$ can unconditionally pay $S$):
\AR{
\mu(S,0) =
g(S) =
L_{21}
\dfrac
{L_{1n}+S}
{L_{1n}+S+L_{12}}
&\geq&
L_{21}
\dfrac
{L_{1n}}
{L_{1n}+L_{12}}
= g(0)
}

\AR{
g'(S) 
&=&
\dfrac
{L_{21}L_{12}}
{(L_{1n}+S+L_{12})^2}
}



\AR{
\mu(S,S) = 
f(S) =
\beta
(S + L_{21}) 
\dfrac
{L_{1n}+S}
{L_{1n}+S+L_{12}}
- S
&\geq&
\beta
L_{21}
\dfrac
{L_{1n}}
{L_{1n}+L_{12}}
= f(0)
}
if $L_{21}=0$ only solution is $S=0$. 

- generalize with assets for 1

- verify that the above is the largest clearing vector

- compute the max of $f$

- add the beta coefficient and see if it still works

- refaire le calcul avec une deformation assymétrique ie $S_{1n}$, $S_{n1}$
différentes; sous les contraintes que $1$ reste insolvable, $n$ reste solvable.

- considerer l'autre magouille de fission, où une banque se scinde en 2 banques (par exemple en mettant tous ses actifs, assets et créances, dans une fille, et toutes les dettes dans l'autre)

- refaisons le calcul en mettant $S$ uniquement du coté de $1$


$f(S) =
\dfrac
{SL_{1n}+SL_{21}+S^{2}+L_{1n}L_{21} - S(L_{1n}+S+L_{12})}
{L_{1n}+S+L_{12}}
$

$f(S) =
\dfrac
{SL_{21}+L_{1n}L_{21} - SL_{12}}
{L_{1n}+S+L_{12}}
$

$f(S) = 
\dfrac
{S(L_{21} - L_{12})+L_{1n}L_{21}}
{L_{1n}+S+L_{12}}
$

Donc 


$f'(S) = 
\dfrac
{(L_{21} - L_{12})(L_{1n}+S+L_{12}) - (S(L_{21} - L_{12})+L_{1n}L_{21})}
{(L_{1n}+S+L_{12})^{2}}
$

$f'(S) = 
\dfrac
{L_{21}L_{1n} - L_{12}L_{1n} + L_{21}L_{12} - L_{12}^{2} - L_{1n}L_{21}}
{(L_{1n}+S+L_{12})^{2}}
$

$f'(S) = 
\dfrac
{(L_{21} - L_{1n} - L_{12})L_{12}}
{(L_{1n}+S+L_{12})^{2}}
$


Donc $f'(S)$ est négative dès que $L_{21} <  L_{1n} + L_{12}$,
cad dès que ce que reçoit le noeud 1 est $<$ ne suffit pas à 
payer ce qu'il doit (ce que nous avons supposé).

donc sous ces hypothèses $f(S)$ est décroissante par conséquent le maximum de $f(S)$ est en $S = 0$ et donc $n$ réalise le maximum possible
$\dfrac
{L_{1n}L_{21}}
{L_{1n}+L_{12}}
$
pour $S=0$ donc mort de la magouille!


\subsection{dérivées}

En cas de prêts assymétriques, soient: 
$S_{1}$ la dette supplémentaire de $1$ envers $n$, et 
$S_{2}$ la dette de $n$ envers $1$.

\AR{
\mu(S_{1},S_{2}) = 
\beta
(S_{2} + L_{21}) 
\dfrac
{L_{1n}+S_{1}}
{L_{1n}+S_{1}+L_{12}}
- S_{2}
\to_{S_1\to\infty}
\beta
(S_{2} + L_{21})- S_{2}
}
si $S_1=S_2$ on a ajouté une paire de créances symétriques ce qui ne changent pas les assets; si les deux dettes sont différentes il vaut mieux utiliser l'équation ci-dessous qui prend en compte la différence:
\AR{
\mu_1(S_{1},S_{2}) = 
(\alpha (S_{1}-S_2) + 
\beta
(S_{2} + L_{21}))
\dfrac
{L_{1n}+S_{1}}
{L_{1n}+S_{1}+L_{12}}
- S_{2} - S_{1}
}
pour expliquer l'asset $(S_{1}-S_2)$ ci-dessus
prenons un exemple
\AR{
1: S_1=+50 & n: -50 + d_{1n}(50)\\
1: -20 + d_{n1}(20)& n: S_2=+20
\\\hline\\
1: 
S_1-S_2=+30 
+ d_{n1}(20)&
n:
S_2-S_1=-30 + d_{1n}(50)
}
donc l'asset en compatibilité locale est bien $S_1-S_2$ du point de vue $1$.


$\triangleright$ 
On veut montrer que le noeud $n$ a intérêt à prêter au noeud $1$.

On définit un jeu à 3 joueurs dans cet exemple.
\begin{definition}\label{jeu3}
Le payoff est déterminé par le (plus grand) vecteur de clearing.
Les coups sont 1) ajouter/retrancher de l'asset,
2) preter/emprunter,
3) rembourser (avant le clearing).
\end{definition}

Concentrons nous sur le deuxieme type de coup.
La formule (2) nous dit le profit du joueur $n$:

+ refaire le calcul d'optimisation avec (2)

+ chercher l'equilibre de Nash avec des coups de type 2) pour commencer

+ analyser le cas où l'ajout de $S_1$ conduit à la faillite de 2; ceci devrait avoir un impact sur la profitabilité de $n$.

\AR{
\dfrac
{\partial \mu}
{\partial S_{1}} =
\beta
(S_{2} + L_{21}) 
\dfrac
{(L_{1n}+S_{1}+L_{12}) - (L_{1n}+S_{1})}
{(L_{1n}+S_{1}+L_{12})^{2}}
\\
\dfrac
{\partial \mu}
{\partial S_{1}} =
\beta
(S_{2} + L_{21}) 
\dfrac
{L_{12}}
{(L_{1n}+S_{1}+L_{12})^{2}}
>
0
}
Donc pour $S_{2}$ fixé $\mu$ est strictement croissante.


\AR{
\dfrac
{\partial \mu}
{\partial S_{2}} =
\beta
\dfrac
{(L_{1n}+S_{1})}
{(L_{1n}+S_{1}+L_{12})}
- 1
< 
0
}
Pour $S_{1}$ fixé $\mu$ est strictement décroissante.

\AR{
\dfrac
{\partial^{2} \mu}
{\partial S_{1} \partial S_{2}} =
\dfrac
{\partial^{2} \mu}
{\partial S_{2} \partial S_{1}} =
\beta
\dfrac
{L_{12}}
{(L_{1n}+S_{1}+L_{12})^{2}}
>
0
}

\AR{
\dfrac
{\partial^{2} \mu}
{\partial S_{1}^{2}} =
\beta
(S_{2} + L_{21})L_{12}
\dfrac
{-2}
{(L_{1n}+S_{1}+L_{12})^{3}}
<
0
}
Donc plus $S_{1}$ augmente et plus la valeur de $\mu$ tend vers une constante.

\AR{
\dfrac
{\partial^{2} \mu}
{\partial S_{2}^{2}} =
0
}

De plus en raisons des hypothèses : $S_{1}; S_{2} \geq 0$, $L_{12} > 0$ (sinon la magouille n'a aucune raison d'être) on a :
\AR{
\nabla \mu(S_{1}; S_{2})
\neq
\vec{0}
}
Par conséquent $\mu$ ne possède pas de points critiques et n'a donc pas de maximum ou de minimum locaux. 


On pose comme contraintes :

- 1 doit rester insolvable donc $L_{21} + S_{2} < L_{12} + L_{1n} + S_{1}$

- n reste solvable donc $S_{1} \geq S_{2}$

\AR{
\mu(S_{1},S_{2}) = 
\beta
(S_{2} + L_{21}) 
\dfrac
{L_{1n}+S_{1}}
{L_{1n}+S_{1}+L_{12}}
- S_{2}
\\
\dfrac
{\partial \mu}
{\partial S_{1}} 
>
0\\
\dfrac
{\partial \mu}
{\partial S_{2}} 
<
0
} 

donc $\mu_{S_{2}}(S_{1})$ est strictement croissante et $\mu_{S_{1}}(S_{2})$ est strictement décroissante par conséquent les contraintes n'en sont pas vraiment.


$\triangleright$ 
Dans le cas $S_1\to\infty$ (magouille max) et $S_2=0$, on voit que
le profit de $n$ tend vers $\beta L_{21}$; donc tout l'argent $L_{21}$ (du par 2 à 1) est détourné par $n$ - alors que cet argent est partagé avec $2$ dans la configuration initiale.

attention dans la magouille $(S_1,0)$ à ne pas mettre 2 en faillite:
ce qui pourrait se produire si la magouille $1n$ 
lui dérive trop de l'argent qui devait lui revenir pour $S_1=0$.

on peut aussi imaginer comme $L_{2n}=0$
que $n$ vient attaquer $2$ en prétant à $1$.

pour un noeud insolvable il y a deux attaques - 1) retirer ses assets, 2) contracter une dette (voir ci-dessus); une contre-mesure pour 2) serait de forcer la fusion (et annuler ainsi la dette); on pourrait également forcer n à emprunter à 2.

le jeu de magouille est compliqué du fait que: 1) il ne faut pas mettre 2 en faillite, 2) que font les autres joueurs (ils peuvent aussi se mettre à magouiller)

$\triangleright$ 
le profit pour le noeud $n$ s'écrit en supposant que 2 soit solvable et 1 insolvable ($n\to^{S_1} 1$, $1\to^{S_2} n$, et on considère que le montant du prêt s'ajoute aux actifs du prêteur):
\AR{
\mu_1(S_{1},S_{2}) = 
(\alpha (S_{1}-S_2) + 
\beta
(S_{2} + L_{21}))
\dfrac
{L_{1n}+S_{1}}
{L_{1n}+S_{1}+L_{12}}
- S_{2} - S_{1}
\\
\mu_1(S_{1},S_{2}) = 
\alpha (S_{1}-S_2) 
\dfrac
{L_{1n}+S_{1}}
{L_{1n}+S_{1}+L_{12}}
 + 
\beta
(S_{2} + L_{21})
\dfrac
{L_{1n}+S_{1}}
{L_{1n}+S_{1}+L_{12}}
- S_{2} - S_{1}
\\
\mu_1(S_{1},S_{2}) = 
\alpha 
\dfrac
{S_{1}L_{1n}+S_{1}^2 - S_{2} L_{1n} - S_{2}S_{1}}
{L_{1n}+S_{1}+L_{12}}
 + 
\beta
(S_{2} + L_{21})
\dfrac
{L_{1n}+S_{1}}
{L_{1n}+S_{1}+L_{12}}
- S_{2} - S_{1}
\\
\dfrac
{\partial \mu_1}
{\partial S_{1}} =
\alpha
\dfrac
{(2S_{1} + L_{1n} - S_{2})(L_{1n}+S_{1}+L_{12}) - (L_{1n}S_{1}+S_{1}^2 -L_{1n}S_{2} - S_{1}S_{2})}
{(L_{1n}+S_{1}+L_{12})^2}
+
\beta
(S_{2} + L_{21}) 
\dfrac
{L_{12}}
{(L_{1n}+S_{1}+L_{12})^{2}}
- 1
}
Or 
\AR{
(2S_{1} + L_{1n} - S_{2})(L_{1n}+S_{1}+L_{12}) = 
3S_{1}L_{1n} + 2S_{1}^2 + 2S_{1}L_{12} +L_{1n}^2 +L_{1n}L_{12} - S_{2}L_{1n} - S_{2}S_{1} - S_{2}L_{12}
}
donc 
\AR{
(2S_{1} + L_{1n} - S_{2})(L_{1n}+S_{1}+L_{12}) - (L_{1n}S_{1}+S_{1}^2 -L_{1n}S_{2} - S_{1}S_{2})
\\
=S_{1}^2 + L_{1n}^2 + 2S_{1}L_{1n} + 2S_{1}L_{12} + L_{1n}L_{12} - S_{2}L_{12} 
\\
=S_{1}^2 + L_{1n}^2 + 2S_{1}L_{1n} + 2S_{1}L_{12} + 2L_{1n}L_{12} + L_{12}^2  - L_{12}^2 - L_{1n}L_{12} - S_{2}L_{12}
\\
=(L_{1n}+S_{1}+L_{12})^2 - L_{12}^2 - L_{1n}L_{12} - S_{2}L_{12}
\\
donc
\\ 
\dfrac
{\partial \mu_1}
{\partial S_{1}} =
\alpha
\dfrac
{(L_{1n}+S_{1}+L_{12})^2 - L_{12}^2 - L_{1n}L_{12} - S_{2}L_{12}}
{(L_{1n}+S_{1}+L_{12})^2}
+
\beta
(S_{2} + L_{21}) 
\dfrac
{L_{12}}
{(L_{1n}+S_{1}+L_{12})^{2}}
- 1
\\
\dfrac
{\partial \mu_1}
{\partial S_{1}} =
\alpha
(1 -
\dfrac
{L_{12}^2 + L_{1n}L_{12} + S_{2}L_{12}}
{(L_{1n}+S_{1}+L_{12})^2})
+
\beta
(S_{2} + L_{21}) 
\dfrac
{L_{12}}
{(L_{1n}+S_{1}+L_{12})^{2}}
- 1
}


Si 
$1 - \dfrac{L_{12}^2 + L_{1n}L_{12} + S_{2}L_{12}}{(L_{1n}+S_{1}+L_{12})^2} > 0$ (*)
 alors soient $\alpha = \beta = 1$.
 
On a par hypothèse  : $L_{21} < L_{1n} + L_{12}$ 
\AR{
L_{21} - L_{1n} - L_{12} 
<
0
\\
L_{21}L_{12} - L_{1n}L_{12} - L_{12}^2
<
0 
\\
(L_{1n}+S_{1}+L_{12})^{2} + L_{21}L_{12} - L_{1n}L_{12} - L_{12}^2 - S_{2}L_{12} + S_{2}L_{12}
< 
(L_{1n}+S_{1}+L_{12})^{2}
\\
1 - 
\dfrac
{L_{12}^2 + L_{1n}L_{12} + S_{2}L_{12}}
{(L_{1n}+S_{1}+L_{12})^2})
+
(S_{2} + L_{21}) 
\dfrac
{L_{12}}
{(L_{1n}+S_{1}+L_{12})^{2}}
<
1
}
Donc pour (*), $\alpha, \beta = 1$ étant le maximum on a : 
$\dfrac {\partial \mu_1} {\partial S_{1}} < 0$
\\

Si 
$1 - \dfrac{L_{12}^2 + L_{1n}L_{12} + S_{2}L_{12}}{(L_{1n}+S_{1}+L_{12})^2} \leq 0$ alors 
\AR{
L_{12}^2 + L_{1n}L_{12} + S_{2}L_{12}
\geq
(L_{1n}+S_{1}+L_{12})^2
\\
S_{2}L_{12} 
\geq
S_{1}^2 + L_{1n}^2 + 2S_{1}L_{1n} + 2S_{1}L_{12} + L_{1n}L_{12}
\\
S_{2} 
\geq
\dfrac{S_{1}^2 + L_{1n}^2 + 2S_{1}L_{1n}}
{L_{12}}
+
2S_{1} + L_{1n}
\\
S_{2} 
\geq
\dfrac{(S_{1} + L_{1n})^2}
{L_{12}}
+
2S_{1} + L_{1n}
}
Donc $S_{2}$ devrait être très supérieur à $S_{1}$ ce qui n'est pas possible puisque par hypothèse on a $S_{2} \leq S_{1}$.

de plus
\AR{
\dfrac
{\partial \mu_1}
{\partial S_{2}} =
(\beta - \alpha)
\dfrac
{(L_{1n}+S_{1})}
{L_{1n}+S_{1}+L_{12}}
- 1 
\leq 
0
\\
\dfrac
{\partial^2 \mu_1}
{\partial S_{1}S_{2}} =
(\beta - \alpha)
\dfrac
{L_{12}}
{(L_{1n}+S_{1}+L_{12})^2}
\\
\dfrac
{\partial^2 \mu_1}
{\partial S_{1}^2} =
2(
\alpha 
\dfrac
{(L_{1n}L_{12}+S_{2}L_{12}+ L_{12}^2)}
{(L_{1n}+S_{1}+L_{12})^3}
-
\beta
\dfrac
{(S_{2}+L_{21})L_{12}}
{(L_{1n}+S_{1}+L_{12})^3})
}


Il n'y a pas de points critiques, les deux dérivées sont négatives pour les contraintes imposées ($S_{1} - S_{2} > 0$ et 1 reste insolvable)  donc la maximum est atteint pour $S_{1}=S_{2}=0$.

\section{Standstill: le jeu à trois joueurs}
\underline{Modélisation d'un prêt}: supposons que $i$ prete $x$ à $j$,
la maj correspondante est la suivante:
\AR{
L_{ji} \leftarrow x + L_{ji}
\\
e_{j} \leftarrow x + e_{j}
\\
e_{i} \leftarrow -x + e_{i}
}
sous la contrainte $x\leq e_i$

\underline{Modélisation d'un remboursement}: supposons que $j$ rembourse $y$ à $i$,
la maj correspondante est la suivante:
\AR{
L_{ji} \leftarrow -y + L_{ji}
\\
e_{j} \leftarrow -y + e_{j}
\\
e_{i} \leftarrow y + e_{i}
}
sous la contrainte $y\leq e_j$

On remarque que: 
1) les deux transitions sont symétriques: $x=-y$,
cad un emprunt negatif = remboursement;
2) l'asset total reste constant (qq part la richesse totale est invariante
et le $L$ dit comment elle est répartie).

La richesse totale ne varie que via l'alea des assets (eg -actions sur marché d'actions)

\begin{theorem}[Standstill theorem]
Dans le jeu fermé~\ref{jeu3} (sans le coup d'entrée/sortie d'asset),
sous les hypothèses idoines de solvabilité, 
la réponse de Nash du joueur $n$ est de
rembourser $\min(L_{n1},e_n)$ à $1$;
et $1$ est indifférent et peut rembourser $\min(L_{1n},e_1)$ par $1$.
\end{theorem}

Conclusion : dans le jeu collusif,
$n$ va partager le gain avec $1$ et les deux joueurs vont rembourser le max possible! De ce fait le joueur $2$ sera lésé 
par violation du principe de non-prioritisation entre dettes
(viol du principe de fongibilité de la dette).

NB: les banques insolvables ont le plus de pouvoir de bargain!

+ valider les coups et la formule $\mu_1$

+ vérifier le thm

+ vérifier que le joueur 2 est bien lésé (perd du payoff)

+ généraliser à un jeu initial qcq à 3 joueurs? determiner dans quels régimes de solvabilité il faut un `standstill'.

+ voire à $n$ joueurs

+ refaire la fonction $\mu_1$ avec 1 qui possède des assets et avec $L_{n1} >0$

On note que la stratégie du joueur $i$ dans le jeu ci-dessus est un simplex:
\AR{
- e_i &\leq& 
\sum_{j\neq i} y_{ij} 
&&\emph{remboursements bornés inférieurement}
\\
% mq Nash eq = clearing?
y_{ij} &\leq& 
e_i
&&\emph{emprunts bornés supérieurement}
}
les coups d'emprunt sont en fait des coups du préteur.

$\partial_{S_1}\mu_1\leq 0$ en $0$,
ce qui veut dire que c'est dans l'intéret du joueur $n$ de rembourser $1$!




\section{readings}

Tim Roughgarden - formule d'allocation de reward sur bitcoin: \url{https://www.youtube.com/watch?v=thyoplWeX1o&t=1s}

\end{document}
